\section*{摘要}

本文报告 Yat-CASched 的设计、实现与实测。Yat-CASched 是一个内置于 Linux 6.8 的缓存感知调度类（策略号 8）：它借鉴缓存重用距离与 Cache Recency Profile（CRP），用作业的历史执行区间估计"把任务放到某个核能省多少时间"，再按最大加速比（MSF）与最小缓存冲击（LCIF）在线选核。本文关心的不只是"设计一个启发式"，而是让它在真实内核里可跑、可测、可归因、可核对。

实现部分给出调度类的数据结构、作业生命周期、唤醒路径、可观测性与亲和性语义，并固化成 8 个可独立应用的补丁：按序应用后与工作树逐字节一致，能编进 \texttt{bzImage} 并在 QEMU 中稳定启动。\texttt{lockdep} 在过程中实证了一处锁反序——\texttt{task\_dead} 在开中断上下文销毁作业，与总是关中断取锁的 balance 路径构成 ABBA；把销毁段放进 \texttt{local\_irq\_save()} 区间后复跑，同类告警计数归零，15 次测量共 30 次独立启动，不再出现告警或 BUG。

实测在 QEMU（4 vCPU，TCG，8 个任务、共 $1.6\times10^7$ 个工作单元）上进行。实验设计上用三组配置做消融：主口径跑完整链路，单核 pin 使 \texttt{nr\_cpus\_allowed==1}、调度类回调一次都不进（\texttt{select\_calls}=0），无唤醒组只剩 \texttt{WF\_FORK}。这三格跑的是同一个二进制、同一份负载、同样的完成量，因此它们的差值是可解释的：完整链路 1.81 倍（全距 $[1.76,2.13]$），只有接入代价时 1.72 倍，无唤醒时 0.84 倍（全距跨 1）。相减得到\textbf{AJLR 选核的边际代价约 5\%}（优化前约 18\%），主体开销被定位到"任务待在这个调度类里"的接入路径，而不是算法本身。

计数器给出决策链路的全貌：512 次唤醒触发 467 次选核（91.2\%），295 次由"目标核空闲"直接复用（63.2\%），172 次做出完整决策并全部落地，接管率 $906/924=98.1\%$、\texttt{place\_miss}=0、\texttt{pick\_null\_desync}=0，平均可运行核数 4.0。计数器同时是优化的入口：\texttt{balance\_runs}=757 而唤醒仅 512 次，据此实现"就绪代际门禁"，使 98.6\% 的 balance 调用成为不取锁的快速路径（\texttt{balance\_runs} $757\to9$、\texttt{balance\_noop}=621），主口径中位数从 2.00 降到 1.81，选核边际代价从约 18\% 压到约 5\%。这次优化的判据是确定性计数器而不是墙钟，因此改完即可判定对错。

本文同时如实记录被推翻的早期结论：项目前期用 Java 模拟器得到的"11.66\% Makespan 改进、30.21 个百分点命中率提升"来自两套模型口径，不可复现；复核后该模拟器已从证据链移除，仅作为设计期校验记录保留在仓库文档中。全文 66 项数字由仓库脚本从归档 CSV 重算，可由 \texttt{make verify-all} 一条命令复核。

\textbf{关键词：}多核调度；缓存感知；DAG 任务；Linux 调度类；可复现测量；QEMU